% Created 2026-08-31 mån 11:09
% Intended LaTeX compiler: xelatex
\documentclass[11pt]{article}

\usepackage{amsmath}
\usepackage{fontspec}
\usepackage{graphicx}
\usepackage{longtable}
\usepackage{wrapfig}
\usepackage{rotating}
\usepackage[normalem]{ulem}
\usepackage{capt-of}
\usepackage{hyperref}
\usepackage[margin=3cm]{geometry}
\author{Felix Loyau-Kahn}
\date{\today}
\title{Assignment, WASP Software Engineering Course Module 2026}
\hypersetup{
 pdfauthor={Felix Loyau-Kahn},
 pdftitle={Assignment, WASP Software Engineering Course Module 2026},
 pdfkeywords={},
 pdfsubject={},
 pdfcreator={},
 pdflang={English}}
\begin{document}

\maketitle
\section{Introduction}
\label{sec:org6a43f5d}
My research focuses on static analysis, particularly Reference Attributed
Grammars (RAGs). RAGs offer a declarative framework for compiler construction in
which attributes are associated with nodes in the Abstract Syntax Tree (AST). By
expressing static analyses and semantic behavior through these attributes, RAGs
enable the development of compilers that are modular, extensible, and
maintainable.

This thesis explores novel strategies for systematically and verifiably
identifying the causes of false-positive and false-negative bug reports in
static checkers. The research focuses on the generation and use of contestable
traces to improve the precision, refinement, and configurability of static
analyses. A contestable trace is a structure that represents an explanation of a
result made by a static analysis, designed so that each step can be inspected,
verified, refined, and contested to identify sources of imprecision.

The work will be carried out within the declarative meta-compilation framework
JastAdd and the ExtendJ Java compiler. Different aspects of contestable trace
generation and analysis refinement will be explored through the implementation
of software artifacts, experimental evaluation, and publication of both
scientific results and tools, with a focus on practicality, automation, and
scalability.
\section{Lecture principles}
\label{sec:orge365ff2}
\subsection{Design By Contract}
\label{sec:org25cd4e6}
In the course, Design By Contract is defined as: “If all the routine’s
preconditions are met by the caller, the routine shall guarantee that all
postconditions and invariants will be true when it completes."

You then state that if the language doesn’t support it, we should approximate
with assertions. I think one of the goals of a compiler based on RAGs is to
fulfill that requirement at a language definition level, providing a framework
to express semantic rules and static analyses, acting like contracts which rely
on assumptions, preventing runtime bugs by catching them at compilation time,
verifying that contracts are respected before execution. That follows the “Crash
early” and “Crash, don’t trash”.
\subsection{Verification and Validation}
\label{sec:org5a9bb06}
V\&V also relates with my work: verification is central to static analysis, as it
strongly participates in checking whether we are building the product right, and
validation is more indirectly related, because better explanation can help
developers determine good specification, checking whether we are building the
right product.

By defining semantic rules that a language should follow, and static analysis
that verifies the implementation of software to ensure they follow the language
specification. This belongs to the first of the three groups you state: static
techniques, dynamic techniques, and process evaluation techniques.
\section{Guest-lecture principles}
\label{sec:org5d31464}
\subsection{Requirements Engineering}
\label{sec:orgc5576f5}
Requirements Engineering (RE) is related to my work: we specify semantic rules
that a language must fulfil, and RE defines a prescriptive statement of intent
that describes an abstract property (requirements specification) that a system
must fulfill.

RE specifies what a software should do, while semantic rules specify what valid
programs are allowed to do. In both cases, it's the same motivation: the cost of
defect removal; the later a defect is detected, the more it will cost to remove
it. It is better to catch it in an early phase or define it as a requirement
from the start.
\subsection{Behavioral Software Engineering}
\label{sec:org71b7d0f}
As defined in the slides: BSE is "applying psychological and social principles
to address issues affecting software development."

This is of interest in my research, developers make predictable mistakes, static
analysis is designed to compensate for these human mistakes like division by
zero, null pointer dereferences, and other developer mistakes. We want to ensure
that many people can work together on a project while ensuring the global
quality of the produced implementation stays high, preventing bugs from going
into production, which could be a potential risk. My work of making contestable
traces makes warnings easier to understand, with better explanations and
influences developer behaviour.
\section{Data scientists, software engineers, and AI engineers}
\label{sec:org290a74b}
Data scientists come from a scientific background and mainly does statistics
and machine-learning algorithms, their main activities are gathering and
cleaning data and developing models, which is a model-centric point of
view. Software engineers focus on delivering software products and has a
system-wide point of view to meet the users and clients requirements, designing
the architecture of a system, implementing, testing, and deploying it, analyzing
risks, and designing safety and security mechanisms into the system, bringing it
to deployment in production-like expectations.

This idea of a system-wide view and model-centric view applies well to my
research area of compilation and static analysis: ExtendJ is a large software
system, and my research may introduce AI-assisted analyses, which need to
integrate into deterministic compiler infrastructure, therefore requiring
system-wide perspective, with the AI part being considered as black-box
functions that could return incorrect results. Our domain of expertise is far
from that of data scientists, being more centered on producing a well-structured
codebase. If AI-assisted static analyses become part of future compiler,
developers will need explanation mechanisms to verify the AI-generated results
before trusting them in production, which should be part of my contestable traces
work.

After reading these chapters, I think data scientists and software engineers
will remain distinct roles. Their skills have little overlap, and they need to
work together to build production ML systems. Specialization remains valuable in
my opinion to keep both system-wide view and model-centric view consistent and
well produced.

But it will become more important that both have some knowledge of the other's
domain: from the data scientists perspective, to meet the needs, requirements,
and hardware constraints of the global project; and from the software engineer's
perspective, to incorporate the models accordingly into a large-scale project
that can be deployed in production with a good level of safety and security.

Some new concepts and technologies tend to change the boundaries between these
roles: foundation models reduce the amount of task-specific model development,
AI coding assistants reduce the barrier for software engineers to use ML,
Retrieval-Augmented Generation (RAG) systems used for domain-specific tasks like
either data or software engineering, and finally, prompt pipelines with agentic
workflows can assist any role and code entire applications with minimal human
intervention, ML-enabled projects included.

The term "AI Engineering" highlights a stronger focus on engineering in
machine-learning projects. It mostly refers to building automated pipelines,
deploying models, MLOps, and it tends to promote system-wide rather than
model-centric points of view, which is becoming increasingly necessary. It is at
the center of these two chapters.
\section{Paper analysis}
\label{sec:org8d4bc6c}
\subsection{First paper: An Exploratory Study of V-Model in Building ML-Enabled Software: A Systems Engineering Perspective}
\label{sec:org41a72e7}
The first paper I chose is written by Jie JW Wu, published in the CAIN 2024
conference: Proceedings of the IEEE/ACM 3rd International Conference on AI
Engineering - Software Engineering for AI.
\subsubsection{Core ideas and their SE importance}
\label{sec:orgd2f0109}
The paper continues the discussion from the first two chapters of “Machine
Learning in Production”, they also observe the challenges of the development
process for ML-enabled software in large-scale real-world applications due to
data quality, technical debt, testing, security \& privacy, safety, and
interdisciplinary collaborations.

They apply a systems engineering lens to investigate the use of the V-Model to
address these challenges, and by establishing several propositions to manage
interdisciplinary collaborations when building products with ML components.

This paper fits very well into the subject of engineering of AI systems and
takes a closer look at the application of a specific software engineering
methodology to find its pros and cons. It opens up future research directions
such as system decomposition, clear system boundaries, and consistency of
Validation \& Verification (V\&V) for building ML-enabled systems.
\subsubsection{Relation to my research}
\label{sec:org26e6a4e}
In my research, while working on a declarative compiler with the aim of semantic
verification and static analysis, we try to keep our software as deterministic
as possible. The subject of building AI-assisted RAG-based declarative compilers
is still an emerging research area.

We are going to try and use AI components in very specific parts of our
compiler. In my research project, for example, to find counterexamples that
explain false positives and negatives in some static analysis, which is normally
NP-hard, but could potentially be solved more efficiently using well-trained
AI. The solutions returned by the AI could easily be verified in linear
complexity.

The integration of such ML components would have to be well isolated from the
rest of the compiler, and since no ML component has been introduced yet, we
still don't have any workflow on how to incorporate it into our software while
keeping good isolation guarantees. This V-model could be an interesting solution
on how to work with ML components for our future work.
\subsubsection{Integration into a larger AI-intensive project}
\label{sec:orgf78ca7f}
Let's imagine a software project that reads and classifies mail and bank checks
automatically for a given list of clients. As a start, we build software that
reads through some PDF and text files and does some pattern matching or fuzzy
matching with a list of clients to sort each mail to the corresponding
client. But then we want to introduce the classification of scanned mail, and to
do so we have to perform optical character recognition (OCR) on the PDF using
AI. Also, if we want to scan, for example, bank checks to automatically put all
the information in a spreadsheet, which contains handwriting, we would also need
AI to extract this handwritten text. So some AI components have to be
introduced.

The difficulties introduced by adding ML components will mainly come from the
collaboration of software and ML engineers, following their modified V-model
lifecycle would help this collaboration by securing the process and enabling
production-ready software.
\subsection{Second paper: Less Is More: Measuring How LLM Involvement affects Chatbot Accuracy in Static Analysis}
\label{sec:org2090275}
The second paper I chose was written by Krishna Narasimhan in 2026. It addresses
several subjects: LLMs for software engineering, AI-assisted development,
LLM/RAG/prompt evaluation, and Agentic software engineering.
\subsubsection{Core ideas and their SE importance}
\label{sec:org48e9e76}
This paper studies how LLMs should interact with a static analysis framework to
answer code-analysis questions. The authors first introduce LLM involvement as a
design variable and compare three architectures along a spectrum of LLM
autonomy. The first translates a natural language description of a code analysis
task to CPGQL queries directly, the second transforms natural language to an
intermediate representation in a JSON table and then passes that table through a
deterministic algorithm to produce the static analyzer queries, and finally, the
third gives full control to the LLM by letting it invoke the tool and pass the
arguments itself. We have three levels of LLM involvement, which are compared in
terms of correctness on known examples, large-scale projects, and LLM token
usage.
\subsubsection{Relation to my research}
\label{sec:org85c6d31}
This paper is closely related to my research. We want to introduce LLMs for very
specific tasks in our compiler, but we want to keep as much determinism as
possible. This paper introduces three levels of LLM involvement in static
analysis and shows that they are not trained on enough static-nalysis data to
produce a DSL with correct semantics, and that it is a better choice to produce
structured output which is then passed to deterministic algorithms to maximize
correctness and optimize token usage.
\subsubsection{Integration into a larger AI-intensive project}
\label{sec:org376d309}
In our compiler, we have developed a lot of static analysis for Java, and we
believe one of the main reasons that prevent people from using it is the
difficulty of tracing false positives and negatives. We have in mind to try to
use LLMs to produce concrete examples that show these errors to the user. This
is an algorithmically complex problem, but one that could be verified in linear
time after being produced by an LLM.

Since my own WASP research focuses on tracing these false positives, we could
also use LLMs to identify which widening operation introduced the imprecision
causing these false positives and negatives, and compare them with the result of
a deterministic algorithm.
\subsubsection{Adaptation of your research}
\label{sec:org51ac938}
The use of LLMs in a compiler is still uncertain, this paper has introduced ways
to differentiate and compare LLM usage in three different ways. When we
introduce AI into our project, we will be able to test different solutions,
compare them, and choose accordingly depending on our needs, also comparing them
with our own deterministic solution.
\section{Evidence, originality, and GenAI-use transparency}
\label{sec:org67dce05}
\begin{itemize}
\item "An Exploratory Study of V-Model in Building ML-Enabled Software: A Systems
Engineering Perspective" is a peer-reviewed paper published in the CAIN
conference proceedings.  \url{https://dl.acm.org/doi/10.1145/3644815.3644951}
\item "Less Is More: Measuring How LLM Involvement affects Chatbot Accuracy in
Static Analysis" is an arXiv paper released in 2026, not yet published in a
peer-reviewed journal or conference proceedings.
\url{https://arxiv.org/abs/2604.21746}
\item "Machine Learning in Production: From Models to Products" is a book that has
been published by MIT Press as open access.  \url{https://mlip-cmu.github.io/book/}
\item The course material.
\item The sources I trust the most are:
\begin{itemize}
\item the first paper because it is peer-reviewed.
\item the second paper because although it is a preprint, it provides experimental
evidence supporting its conclusions.
\item the MIT Press book is also highly trustworthy because it is an academic
textbook.
\end{itemize}
\item I used GenAI for language checking, after I finished the assignment.
\end{itemize}
\section{Conclusion}
\label{sec:orga02dc54}
This course was relevant for my future work which will probably include AI. It
gave me insight into the current practices in AI engineering and how to use it
safely and methodically to benefit fully from its capabilities while keeping
software reliable and safe. It gave me a methodology for integrating AI
components in large-scale projects, while keeping in mind the production
constraints, discussing the many difficulties encountered by people working with
AI in large-scale projects involving interdisciplinary teams.

For my research, after doing this course, I think AI should assist deterministic
compiler analyses rather than replace them.

The course also changed the way I think about AI usage. Before the course, I
viewed AI as an isolated component, whereas I now see it more as a part of a
larger engineering problem to incorporate with a system-wide view. I really
appreciated the importance of verification and production concerns.
\end{document}
